% Part: normal-modal-logic
% Chapter: axioms-systems
% Section: derived-rules

\documentclass[../../../include/open-logic-section]{subfiles}

\begin{document}

\olfileid{nml}{prf}{der}

\olsection{નિષ્પન્ન નિયમો}

!!{derivation}s શોધવી અને લખવી સ્પષ્ટપણે મુશ્કેલ, કંટાળાજનક અને
પુનરાવર્તિત કામ છે. દાખલા તરીકે, ઘણી વાર આપણે $!A \lif !B$માંથી
$\Box !A \lif \Box !B$ મેળવવું હોય છે, એટલે કે નિયમ~\RK વાપરવો હોય
છે. તે માટે પહેલાં \Nec વાપરવું, પછી~\Ax{K}નો યોગ્ય દૃષ્ટાંત લખવો
અને છેલ્લે~\MP વાપરવું પડે. $!A \lif
!B$ અને $!B \lif !C$માંથી $!A \lif !C$ મેળવવા માટે આ લાંબો
પુનરુક્તિ-દૃષ્ટાંત લખવો પડે:
\[
(!A \lif !B) \lif ((!B \lif !C) \lif (!A \lif !C))
\]
અને \MP{} બે વાર વાપરવું પડે. ઘણી વાર કોઈ ઉપ-!!{formula}ને તેની સાથે
સમકક્ષ હોવાનું આપણને જાણીતું હોય એવા સૂત્રથી બદલવું હોય છે; દાખલા
તરીકે, \iftag{prvDiamond}{$\Diamond
  !A$ને $\lnot\Box\lnot !A$થી}{$\Box !A$ને $\lnot\Diamond\lnot !A$થી}, અથવા
$\lnot\lnot !A$ને $!A$થી. તેથી આખી !!{derivation} લખવાને બદલે
વચ્ચેનાં પગલાં !!{derivable} કેમ છે એટલું નોંધવું વધુ અનુકૂળ છે. આ માટે
!!{derivability} વિશે કેટલીક હકીકતો એકત્ર કરીએ.

\begin{prop}
  જો $\Log{K} \Proves !A_1$, \dots, $\Log{K} \Proves !A_n$ હોય અને
  $!A_1$, \dots, $!A_n$માંથી વિધાનાત્મક તર્ક અનુસાર $!B$ ફલિત થતું હોય,
  તો $\Log{K} \Proves !B$.
\end{prop}

\begin{proof}
  જો $!A_1$, \dots,~$!A_n$માંથી વિધાનાત્મક તર્ક અનુસાર $!B$ ફલિત થતું હોય, તો
  \[
  !A_1 \lif (!A_2 \lif \cdots (!A_n \lif !B)\dots)
  \]
  પુનરુક્તિનો દૃષ્ટાંત છે. \MP{}ને $n$ વાર વાપરવાથી~$!B$ની
  !!{derivation} મળે છે.
\end{proof}

આ પ્રસ્તાવના ઉપયોગને આપણે~\PL વડે દર્શાવીશું.

\begin{prop}
  જો $\Log{K} \Proves !A_1 \lif (!A_2 \lif \cdots (!A_{n-1} \lif
  !A_n)\dots)$, તો $\Log{K} \Proves \Box !A_1 \lif (\Box !A_2 \lif
  \cdots (\Box !A_{n-1} \lif \Box !A_n)\dots)$.
\end{prop}

\begin{proof}
  \olref[nor]{prop:rk}ની સાબિતીની જેમ $n$ પર આગમનથી.
\end{proof}

આ પ્રસ્તાવના ઉપયોગને આપણે~\RK વડે દર્શાવીશું. આ પરિણામો વડે
!!{derivability} વિશેના દાવા કેવી રીતે સહેલાઈથી સ્થાપી શકાય તે જોઈએ.

\begin{prop}
  $\Log{K} \Proves (\Box!A \land \Box!B) \lif \Box (!A \land !B)$
\end{prop}

\begin{proof}
  \begin{derivation}
    1. & $\Log{K} \Proves !A \lif (!B \lif (!A \land !B))$ & \Taut \\
    2. & $\Log{K} \Proves \Box!A \lif (\Box !B \lif \Box(!A \land !B)))$ & \RK, 1\\
    3. & $\Log{K} \Proves (\Box!A \land \Box!B) \lif \Box(!A \land !B)$ & \PL, 2
  \end{derivation}
\end{proof}

\begin{prop}\ollabel{prop:rewriting}
  જો $\Log{K} \Proves !A \liff !B$ અને $\Log{K} \Proves
  \Subst{!C}{!A}{q}$ હોય, તો $\Log{K} \Proves \Subst{!C}{!B}{q}$.
\end{prop}\sourcecorrection{OLMD-025}{મૂળ પ્રસ્તાવના અંતિમ પ્રતિસ્થાપનમાં B પહેલાં સૂત્રચિહ્ન ! છૂટી ગયું છે; પૂર્વપક્ષ અને પછીના પ્રશ્નની જેમ અહીં !B મૂક્યું છે.}

\begin{proof}
  અભ્યાસપ્રશ્ન.
\end{proof}

\begin{prob}
  $!C$ની જટિલતા પર આગમન વડે, જો $\Log{K} \Proves !A \liff !B$ હોય,
  તો $\Log{K} \Proves \Subst{!C}{!A}{q} \liff \Subst{!C}{!B}{q}$
  સાબિત કરો અને તેના વડે \olref[nml][prf][der]{prop:rewriting} સાબિત કરો.
\end{prob}

આ પ્રસ્તાવ ખાસ કરીને $\Diamond$ને $\Box$માં (અથવા વિપરીત રીતે)
ફેરવવો હોય, કે !!a{formula}ની અંદર બેવડો નિષેધ દૂર કરવો હોય, ત્યારે
ઉપયોગી છે. આગળ જ્યારે~$!C(!A)$ને !!a{formula}~$!C(!B)$ તરીકે ફરી
લખીશું, ત્યારે \olref{prop:rewriting}ના ઉપયોગને ``$!A$ના સ્થાને $!B$''
વડે દર્શાવીશું. બીજા શબ્દોમાં, ``$!A$ના સ્થાને $!B$'' આનું સંક્ષેપ છે:\sourcecorrection{OLMD-026}{મૂળ સામાન્ય સંક્ષેપમાં ``A for B'' છે, છતાં તેની નીચેની ત્રણ પંક્તિમાં C(A)માંથી C(B) મળે છે અને પછીના ઉદાહરણમાં p જૂના બેવડા નિષેધને બદલે છે; બદલીની દિશા અહીં સ્પષ્ટ કરી છે.}
  \begin{derivation}
    & $\Proves !C(!A)$ \\
    & $\Proves !A \liff !B$\\
    & $\Proves !C(!B)$ & \olref{prop:rewriting} વડે
  \end{derivation}
દાખલા તરીકે:

\begin{prop}
  $\Log{K} \Proves \lnot\Box p \lif \Diamond \lnot p$
\end{prop}

\begin{proof}
  \iftag{prvDiamond}{% Diamond મૂળભૂત છે
  \begin{derivation}
    1. & $\Log{K} \Proves \Diamond \lnot p \liff \lnot\Box\lnot\lnot p$ &
    \Dual\\
    2. & $\Log{K} \Proves \lnot \Box \lnot\lnot p \lif \Diamond \lnot p$ & \PL, 1\\
    3. & $\Log{K} \Proves \lnot\Box p \lif \Diamond \lnot p$ & $\lnot\lnot p$ના સ્થાને $p$\\
  \end{derivation}
  }{% Diamond વ્યાખ્યાયિત છે
  \begin{derivation}
    1. & $\Log{K} \Proves \lnot \Box \lnot\lnot p \lif \Diamond \lnot p$ & \Taut\\
    2. & $\Log{K} \Proves \lnot\Box p \lif \Diamond \lnot p$ & $\lnot\lnot p$ના સ્થાને $p$\\
  \end{derivation}
  પંક્તિ~$1$ પરનું સૂત્ર $\lnot q \lif \lnot q$નો દૃષ્ટાંત છે, જેમાં
  $q$ના સ્થાને $\Box\lnot\lnot p$ મૂક્યું છે. $\Diamond\lnot p$ અને
  $\lnot\Box\lnot\lnot p$ એક જ સૂત્ર છે, કારણ કે $\Diamond$ની
  વ્યાખ્યા~$\lnot\Box\lnot$ છે.}
\end{proof}

ઉપરની !!{derivation}માં છેલ્લું પગલું ``$\lnot\lnot p$ના સ્થાને $p$''
આનો સંક્ષેપ છે:
  \begin{derivation}
    & $\Log{K} \Proves \lnot \Box \lnot\lnot p \lif \Diamond \lnot
    p$\\
    & $\Log{K} \Proves \lnot\lnot p \liff p$ & \Taut\\
    & $\Log{K} \Proves \lnot\Box p \lif \Diamond \lnot p$ & \olref{prop:rewriting} વડે
  \end{derivation}
અહીં \olref{prop:rewriting}ના $!C(q)$, $!A$ અને $!B$ની ભૂમિકામાં અનુક્રમે
$\lnot \Box q \lif \Diamond \lnot p$,
$\lnot\lnot p$ અને~$p$ આવે છે.

જ્યારે !!a{formula}માં ઉપ-!!{formula} $\lnot\Diamond !A$ હોય, ત્યારે
\olref{prop:rewriting} વડે તેની જગ્યાએ $\Box\lnot !A$ મૂકી શકીએ,
કારણ કે $\Log{K} \Proves \lnot\Diamond !A \liff \Box\lnot !A$.
આવા અને સમાન પ્રતિસ્થાપનોને આપણે માત્ર ``$\lnot\Diamond$ના સ્થાને
$\Box\lnot$'' વડે દર્શાવીશું.

નીચેનો પ્રસ્તાવ !!{derivability} વિશેનાં પરિણામો ઢાંચારૂપે સ્થાપવાનું
સમર્થન આપે છે. દાખલા તરીકે, અગાઉનો પ્રસ્તાવ માત્ર
$\Log{K} \Proves \lnot\Box p \lif
\Diamond \lnot p$ જ નહિ, પરંતુ મનસ્વી~$!A$ માટે
$\Log{K} \Proves \lnot\Box !A \lif \Diamond
\lnot !A$ પણ સ્થાપે છે.

\begin{prop}
  જો $!A$ એ $!B$નો પ્રતિસ્થાપન દૃષ્ટાંત હોય અને $\Log{K} \Proves !B$,
  તો $\Log{K} \Proves !A$.
\end{prop}

\begin{proof}
  આખી !!{derivation}માં પ્રતિસ્થાપન કરતાં પણ યોગ્ય !!{derivation}
  મળે છે તે ચકાસવું લાંબું, પરંતુ નિયમિત કામ છે ($!B$ની
  !!{derivation}ની લંબાઈ પર આગમન કરો). ખાસ કરીને, પુનરુક્તિના
  દૃષ્ટાંતોના પ્રતિસ્થાપન દૃષ્ટાંતો પોતે પુનરુક્તિના દૃષ્ટાંતો છે;
  \iftag{prvDiamond}{\Dual{} અને~}{}\Ax{K}ના દૃષ્ટાંતોના પ્રતિસ્થાપન
  દૃષ્ટાંતો પોતે~\iftag{prvDiamond}{\Dual{}
    અને~}{}\Ax{K}ના દૃષ્ટાંતો છે; અને પૂર્વપક્ષ(ઓ) તથા નિષ્કર્ષમાં
  !!{propositional variable}sના સ્થાને !!{formula}s મૂકતાં \MP{} અને~\Nec{}
  નો ઉપયોગ યોગ્ય જ રહે છે.
\end{proof}


\end{document}
